% III. Проектиране на системата.
\section{Проектиране на системата}
\label{sec:design}

Проектирането изхожда от едно изискване: трите неща, които се менят в
експериментите, тоест задачата, метаевристиката и схемата за паралелно изпълнение,
трябва да се менят независимо едно от друго. Ако смяната на схемата за изпълнение
налага промяна в кода на алгоритъма, всяко сравнение между схемите става сравнение и
между две различни реализации на алгоритъма, тоест престава да бъде измерване на
това, което се твърди.

\subsection{Слоеве и отговорности}
\label{subsec:design-layers}

Системата е разделена на четири слоя. Най-долният описва задачата и предоставя
представяне на решение, целева функция и оператор за преминаване към съседно решение.
Над него стои слоят на метаевристиките, който работи само през този договор и не знае
коя конкретна задача решава. Третият слой съдържа схемите за паралелно изпълнение,
които управляват няколко екземпляра на търсенето и обмена между тях. Най-горният слой
е конфигурирането и измерването: избор на инстанция, на алгоритъм, на схема и на
критерий за спиране, и записване на резултатите в машинно четим вид.

Посоката на зависимостите е само надолу и е показана на Фиг.~\ref{fig:layers}
\figref{fig:layers}. Слоят на задачите не знае нищо за метаевристиките, слоят на
метаевристиките не знае нищо за нишките, а слоят на схемите не знае коя задача се
решава. Практическата последица е, че добавянето на четвърта задача не изисква
промяна в нито един от четирите алгоритъма, а добавянето на пета схема за изпълнение
не изисква промяна в нито една задача.

\begin{figure}[H]
\centering
\begin{tikzpicture}[node distance=0.9cm]
  \node[block, minimum width=11cm] (harness) {Слой 4. Измервателна рамка\\
    \code{Experiment}, \code{Trial}, \code{Summary}: инстанция, алгоритъм, схема,
    бюджет, брой повторения, обобщаване};
  \node[block, minimum width=11cm, below=of harness] (schemes) {Слой 3. Схеми за паралелно изпълнение\\
    \code{run\_multistart}, \code{run\_cooperative}, \code{run\_island};
    \code{SharedBest}, \code{Deadline}, \code{WorkerMetrics}};
  \node[block, minimum width=11cm, below=of schemes] (algos) {Слой 2. Метаевристики\\
    \code{Metaheuristic}: \code{step}, \code{best}, \code{inject}, \code{evaluations}\\
    отгряване, табу търсене, генетичен алгоритъм, мравчена колония};
  \node[block, minimum width=11cm, below=of algos] (problem) {Слой 1. Договор за задача\\
    \code{Problem}: \code{evaluate}, \code{delta}, \code{apply}, \code{random\_move},
    \code{construct}, \code{crossover}, \code{construct\_aco}, \code{deposit}\\
    търговски пътник, раница 0/1, оцветяване на граф};
  \draw[line] (harness) -- (schemes);
  \draw[line] (schemes) -- (algos);
  \draw[line] (algos) -- (problem);
  \node[note, right=0.6cm of algos, text width=3.2cm] (n1) {Слой 2 не знае колко нишки има};
  \node[note, right=0.6cm of problem, text width=3.2cm] (n2) {Слой 1 не знае кой алгоритъм го вика};
\end{tikzpicture}
\caption{Слоеве и посока на зависимостите}
\label{fig:layers}
\end{figure}

\subsection{Договор за задача}
\label{subsec:design-problem}

Договорът за задача е минималният набор от операции, които всяка метаевристика
изисква. Той включва тип на решението, оценяване на решението, генериране на начално
решение и обхождане на съседството. Отделно се предвижда операция за
\term{инкрементално оценяване}, тоест изчисляване на промяната в целевата функция при
конкретен ход, без пълно преизчисляване. Тази операция не е удобство, а проектно
решение с пряко следствие върху производителността: при задачата на търговския пътник
пълното оценяване е с линейна сложност спрямо броя на върховете, а промяната при
размяна на две ребра се изчислява с константен брой операции. Измереното отношение
между двете е дадено в Раздел~\ref{sec:results}.

Окончателният списък от операции е фиксиран след реализацията и на трите задачи, а не
след първата, за да не бъде обобщен от един частен случай. Операциите се разделят на
четири групи:

\begin{itemize}[topsep=0pt,itemsep=2pt,leftmargin=*]
\item \textbf{Оценяване.} \code{evaluate} за пълно пресмятане, \code{delta} за
      промяната при ход, \code{apply} за прилагане на хода заедно с обновяване на
      кешираното състояние, \code{feasible} за проверка на допустимост.
\item \textbf{Генериране.} \code{construct} за конструктивна евристика и
      \code{random\_solution} за случайно допустимо решение. Конструктивната евристика
      е рандомизирана: детерминистична такава би дала на всички нишки едно и също
      начално решение и би обезсмислила независимия многократен старт.
\item \textbf{Съседство.} \code{random\_move} и \code{move\_key}, като вторият връща
      атрибута на хода, използван от табу списъка.
\item \textbf{Специфични за алгоритъм.} \code{crossover} за генетичния алгоритъм и
      \code{pheromone\_size}, \code{construct\_aco} и \code{deposit} за мравчената
      колония. Тези операции принадлежат на задачата, а не на алгоритъма, защото
      равномерна рекомбинация би разрушила пермутация, а рекомбинация с подредба е
      безсмислена за двоичен вектор.
\end{itemize}

\subsubsection{Единно представяне и кеширано състояние}
\label{subsubsec:design-representation}

И трите задачи използват едно и също представяне, вектор от цели числа. При задачата
на търговския пътник това е пермутация на градовете, при раницата е индикаторен
вектор, а при оцветяването е цвят за всеки връх. Единното представяне прави хода,
рекомбинацията и миграцията една и съща операция върху паметта навсякъде и позволява
договорът да бъде обикновен абстрактен клас без шаблони.

Само решението обаче не е достатъчно. Инкременталното оценяване при раницата изисква
текущото общо тегло, а генерирането на ход при търговския пътник изисква обратната
пермутация, тоест позицията на всеки град в маршрута. И двете са част от състоянието
на кандидата, а не на задачата, защото задачата се използва едновременно от няколко
нишки и трябва да остане неизменяема. Затова кандидатът е тройка от решение, скалар и
целочислен масив, а тяхното значение се определя от задачата
\tabref{tab:candidate-state}.

\begin{table}[H]
\centering
\caption{Кеширано състояние на кандидата по задачи}
\label{tab:candidate-state}
\begin{adjustbox}{max width=\textwidth}
\begin{tabular}{@{}L{3.6cm}L{3.3cm}L{3.3cm}L{3.4cm}@{}}
\toprule
\textbf{Задача} & \textbf{Решение} & \textbf{Скалар} & \textbf{Масив} \\
\midrule
Търговски пътник & пермутация & не се ползва & обратна пермутация \\
Раница 0/1 & индикатор 0/1 & общо тегло & не се ползва \\
Оцветяване & цвят на връх & не се ползва & не се ползва \\
\bottomrule
\end{tabular}
\end{adjustbox}
\end{table}

\subsection{Договор за метаевристика}
\label{subsec:design-metaheuristic}

Метаевристиката се разглежда като обект със състояние, който изпълнява стъпки, а не
като функция, която връща решение. Причината е в паралелните схеми: за да може една
нишка да получи решение отвън и да продължи от него, търсенето трябва да бъде
прекъсваемо в определени точки, а състоянието му да бъде видимо. Договорът включва
изпълнение на една стъпка, четене на текущото най-добро решение, приемане на решение
отвън и брояч на извършените оценявания.

Една стъпка означава различно нещо при различните алгоритми: един ход при отгряването,
едно поколение при генетичния алгоритъм и едно преминаване на цялата колония при
мравчения алгоритъм. Затова бюджетът никога не се задава в стъпки, а в изтекло време
или в брой оценявания на целевата функция. По същата причина периодът на обмен между
нишките се брои в оценявания, а не в итерации: период, разумен за отгряването, което
изпълнява милиони итерации в секунда, означава нула обмени за мравчената колония,
която изпълнява около сто и петдесет.

\subsection{Схеми за паралелно изпълнение}
\label{subsec:design-parallel}

И трите схеми използват един и същ договор за метаевристика и се различават само по
това какво се обменя и кога \figref{fig:schemes}.

\begin{itemize}[topsep=0pt,itemsep=2pt,leftmargin=*]
\item \textbf{Независим многократен старт.} Всяка нишка изпълнява самостоятелно
      търсене със собствен генератор на псевдослучайни числа. Обмен няма, а
      резултатът е най-доброто от всички решения. Тази схема е и еталонът, спрямо
      който се измерва цената на комуникацията в останалите две.
\item \textbf{Кооперативна схема със споделено най-добро решение.} Нишките четат и
      обновяват обща най-добра стойност. Обновяването е рядко, а четенето е често,
      тоест схемата се проектира около това съотношение: цената се чете от атомарна
      променлива без заключване, а самото решение се защитава с мютекс.
\item \textbf{Островен модел.} Нишките са разделени на подпопулации, които обменят
      решения през пощенски кутии по пръстен. Топологията е пръстен, а не пълен граф,
      защото при пълен граф едно добро решение се разпространява за един обмен и
      схемата се изражда в кооперативната.
\end{itemize}

\begin{figure}[H]
\centering
\begin{tikzpicture}[node distance=0.7cm]
  % multistart
  \node[accent] (m0) {нишка 0};
  \node[accent, below=of m0] (m1) {нишка 1};
  \node[accent, below=of m1] (m2) {нишка 2};
  \node[above=0.2cm of m0] {\small\textbf{многократен старт}};
  % cooperative
  \node[accent, right=3.2cm of m0] (c0) {нишка 0};
  \node[accent, below=of c0] (c1) {нишка 1};
  \node[accent, below=of c1] (c2) {нишка 2};
  \node[danger, right=1.4cm of c1] (best) {споделено\\най-добро};
  \node[above=0.2cm of c0] {\small\textbf{кооперативна}};
  \draw[line, <->] (c0) -- (best);
  \draw[line, <->] (c1) -- (best);
  \draw[line, <->] (c2) -- (best);
  % island
  \node[accent, right=3.6cm of c0.east, xshift=1.6cm] (i0) {остров 0};
  \node[accent, below=of i0] (i1) {остров 1};
  \node[accent, below=of i1] (i2) {остров 2};
  \node[above=0.2cm of i0] {\small\textbf{островен модел}};
  \draw[line] (i0) -- (i1);
  \draw[line] (i1) -- (i2);
  \draw[line] (i2.west) -- ++(-1.0,0) -- ++(0,3.0) -- (i0.west);
\end{tikzpicture}
\caption{Трите схеми за паралелно изпълнение и обменът между работниците}
\label{fig:schemes}
\end{figure}

\subsection{Управление на състоянието и на случайността}
\label{subsec:design-state}

Възпроизводимостта е изискване към проекта, а не пожелание. Всяка нишка получава
собствен генератор на псевдослучайни числа с извлечено от общо начално число подсеме,
така че едно изпълнение да може да бъде повторено точно при същия брой нишки.
Изпълнение с различен брой нишки не е и не може да бъде побитово еднакво, затова
сравненията между конфигурации се правят по разпределението на резултатите от няколко
независими изпълнения, а не по единично изпълнение.

Възпроизводимостта е точна само при бюджет, зададен в брой оценявания: всяка нишка
получава равен дял и спира на него, независимо от скоростта на машината. При бюджет,
зададен във време, броят итерации зависи от натоварването и резултатът не е побитово
повторим. Това е свойство на измерването, а не дефект, и е записано като такова.

\subsection{Точки, в които проектът очаква измерване}
\label{subsec:design-instrumented}

Измервателните точки са фиксирани заедно с методиката от Раздел~\ref{sec:method} и се
събират в структура за всяка нишка поотделно
\tabref{tab:instrumentation}. Структурата е подравнена на 64 байта, тоест на
кеш линия, а дали това подравняване има значение е измерено, а не предположено;
резултатът е в §\ref{subsec:res-falsesharing}.

\begin{table}[H]
\centering
\caption{Величини, записвани за всяка нишка}
\label{tab:instrumentation}
\begin{adjustbox}{max width=\textwidth}
\begin{tabular}{@{}L{4.2cm}L{9.6cm}@{}}
\toprule
\textbf{Величина} & \textbf{Какво отговаря} \\
\midrule
\code{iterations} & колко стъпки е направил работникът \\
\code{evaluations} & колко оценявания на целевата функция, единната мярка за работа \\
\code{improvements} & колко пъти работникът е подобрил собственото си най-добро \\
\code{offers} & колко пъти е публикувал решение към споделено състояние \\
\code{adopted} & колко пъти е приел решение отвън \\
\code{cas\_failures} & неуспешни атомарни обновявания, тоест мярка за състезание \\
\code{sync\_wait\_ns} & време, прекарано вътре в споделеното най-добро или в кутия \\
\bottomrule
\end{tabular}
\end{adjustbox}
\end{table}
